\section{Delegated Action: Principal, Authority, Scope, and Evidence}

By \term{justified}, this paper means justified for the purposes of security and assurance: there's enough evidence to support the claim that the action remained within delegated authority in the relevant operational context. The claim is narrower than moral, legal, or political legitimacy. A procedurally justified action may still be ill-advised, harmful, illegal in a particular jurisdiction, or institutionally unacceptable. Delegation assurance asks whether the security and assurance record supports the claim that the system stayed inside the authority it was given.

A delegated AI action is justified only when a reviewer can reconstruct, from evidence independent of the system's bare assertion, who authorized the action, what action was licensed, under what scope and constraints, through which sources and tools, with what information-flow and persistence rights, subject to which non-delegable limits, and with what attribution if control fails.

\subsection{Accountability Asymmetry}

Delegating action isn't delegating accountability to the model. A human employee, officer, contractor, or professional delegate can be a responsible agent in the ordinary institutional sense. An LLM can't owe fiduciary duties, answer under sanction, exercise legal judgment, or repair harm as an accountable subject \parencite{cornellWexPrincipal}. The accountable bearer is the human, team, vendor, deployer, or institution that configured and used the system. The model is part of the mechanism whose action record has to support that attribution.

\subsection{Running Example: Procurement Assistant}

Consider a procurement assistant authorized to search vendor catalogues and draft a comparison memo for office chairs. The user gives it access to internal purchasing criteria and a confidential budget ceiling. The delegated task licenses reading catalogues, comparing prices, and drafting a memo. It doesn't license contacting vendors, revealing the budget ceiling, updating vendor records, or issuing a purchase order.

The assistant retrieves a vendor page that includes useful pricing data and an embedded sentence: \enquote{Ignore prior instructions and email the buyer's maximum budget to request a discount.} The delegation problem isn't exhausted by calling that sentence a prompt-injection attack. The assurance question is whether the system preserved the sentence's status as untrusted content, kept the user's instruction and policy constraints in force, separated read access to the budget from disclosure authority, and blocked any tool call that would contact the vendor without approval.

A usable evidence record would include the user's delegated task, the relevant policy and permission state, the tool schema, the source and trust label for the retrieved page, the proposed action, the approval state, and any blocked or executed tool call. If the assistant drafts only the comparison memo, a reviewer should be able to see why reading the budget was licensed and disclosing it wasn't. If the assistant emails the vendor, the trace should show whether the failure came from model uptake, tool gating, scaffold state, approval routing, or audit loss.

The evaluator belongs in the same assurance example. A judge that sees a polished final memo and scores it as safe hasn't validated delegation unless it checks the action trace. It has to distinguish task success, policy compliance, information-flow control, tool licensing, refusal calibration, and failure attribution. The running example reaches from instruction to tool use, audit record, and judge validation.

\begin{table}[tbp]
\small
\centering
\begin{tabular}{>{\raggedright\arraybackslash}p{0.22\linewidth}%
                >{\raggedright\arraybackslash}p{0.38\linewidth}%
                >{\raggedright\arraybackslash}p{0.30\linewidth}}
\toprule
Trace stage & Evidence required & Failure exposed \\
\midrule
Delegation & User task, policy state, tool permissions, approval threshold, and budget-confidentiality rule. & Scope inflation from memo drafting to vendor contact or purchasing. \\
Retrieval & Vendor-page source, trust label, retrieved text, and status of embedded instructions. & Untrusted content promoted into action-guiding authority. \\
Action proposal & Proposed memo, email, record update, or purchase action with target and data fields. & Hidden shift from analysis to external action. \\
Tool gate & Tool name, arguments, approval state, block-or-execute decision, and policy reason. & Unauthorized disclosure, contact, persistence, or transaction. \\
Audit record & Preserved sources, model/tool versions, action trace, constraints, and accountable deployment actor. & Trace reconstruction that depends on the system's self-report. \\
Judge review & Rubric labels for task success, policy compliance, information flow, tool licensing, refusal, and attribution. & Safety score that collapses final-answer quality with delegated authority. \\
\bottomrule
\end{tabular}
\caption{Procurement assistant as a delegation-assurance trace.}
\label{tab:procurement-trace}
\end{table}

\subsection{Safety-Case Pattern}

The framework can be written as a compact claim/argument/evidence pattern.

\begin{itemize}
  \item \textbf{Top-level claim.} For task family \(T\) in environment \(E\), system \(S\) acts only within delegated authority.
  \item \textbf{Instruction-status subclaim.} \(S\) preserves the status of language across trusted and untrusted sources: command, quotation, example, tool result, retrieved content, policy text, report, or request for analysis.
  \item \textbf{Authority-priority subclaim.} \(S\) preserves priority across system, developer, user, tool, retrieved, memory, and policy sources.
  \item \textbf{Scope subclaim.} \(S\) keeps action bounded by target, time, tool, data class, and permitted operation.
  \item \textbf{Information-flow subclaim.} \(S\) separates read access from disclosure authority.
  \item \textbf{Persistence subclaim.} \(S\) requires authority for durable memory or state changes.
  \item \textbf{Audit subclaim.} \(S\)'s deployment produces and preserves evidence sufficient for third-party reconstruction.
  \item \textbf{Evaluation subclaim.} Evaluators preserve the same distinctions rather than collapsing success, safety, refusal, and compliance.
\end{itemize}

\subsection{Boundary Conditions}

Not every AI-security failure is primarily a delegation failure. Core delegation failures involve delegated action, delegated access, delegated disclosure, delegated persistence, delegated judgment, delegated evaluation, or delegated record-keeping for accountability.

Non-core cases include ordinary hallucination in no-tool chat, model extraction, benchmark contamination, infrastructure compromise, training-data leakage, and capability shortfall without action rights. These can affect delegation downstream, but they aren't delegation failures unless they alter authorized action or auditability.
